% !TeX spellcheck = de_DE
% Auto-konvertiert aus BA Entwurf.docx. Inhalt sprachlich/fachlich bitte final gegenpruefen.

\chapter{Arbeitsnotizen und Quellenmaterial}
\label{ch:notizen}

\section{Abkürzungen}

\section{Quellen}

\section{Überschriften}

\section{Unterpunkte}

\section{Eigene Aussagen(Temporär bis Quelle gefunden)}

\section{Dateien}

\section{Grundlagen:}

Liu: s.26-31 RTS \& Embedded Systems

Das: s.34-39 Embedded Systems\& MCU

Das: s.47-49 MCU

\section{Benchmarks:}

Das: s.52-X Definitionen für ram, rom, interrupts etc.

\section{Methodik Erweiterung (Noch nicht fertig, aus Zephyr Doku, einige davon wurden verwendet)}

\subsection{Optimierungen}

Zephyr lässt sich hinsichtlich der Leistung, dem Energieverbrauch und dem Speicherbedarf (Footprint) optimieren.

\section{Optimierung des Speicherbedarfs}

\section{Stack Größen}

Die Stäck-Größen der verschiedenen Systemthreads lassen sich für jede Anwendung anpassen. Dabei sind folgende Standardwerte gesetzt:

CONFIG\_ISR\_STACK\_SIZE: 2048

CONFIG\_MAIN\_STACK\_SIZE: 1024

CONFIG\_IDLE\_STACK\_SIZE: 320

CONFIG\_SYSTEM\_WORKQUEUE\_STACK\_SIZE: 1024

CONFIG\_PRIVILEGED\_STACK\_SIZE: 1024

\section{Peripheriegeräte}

Einige Peripheriegeräte sind standardmäßig aktiviert. Diese lassen sich ebenfalls folgendermaßen deaktivieren:

CONFIG\_GPIO=n

CONFIG\_SPI=n

\section{Debug-/ Informationsoptionen}

Die untenstehenden Optionen bieten die Möglichkeit, mehr Informationen über die laufende Anwendung auszugeben und dienen primär als Mittel für Debugging und Fehlerbehandlung (Error-Handling):

CONFIG\_BOOT\_BANNER

CONFIG\_DEBUG

Der Boot Banner ist hierbei standardmäßig aktiviert und kann einige Bytes in Anspruch nehmen.

\section{MPU/ MMU-Unterstützung}

Um weiterhin Speicher zu sparen und die Leistung zu verbessern, kann die MPU/ MMU-Unterstützung deaktiviert werden. Dadurch geht jedoch die erweiterte Stack-Überprüfung und -Unterst+ützun verloren.

\subsection{Optimierungswerkzeuge}

Zephyr verfügt über einige Werkzeuge, welche die Möglichkeit bieten, den Footprint, die Speichernutzung und Datenstrukturen über verschiedene Build-System-Ziele (Targets) analysieren zu können.

\section{Footprint und Speichernutzung}

Für die Ansicht und Analyse von RAM-, ROM-, und der Stack-Nutzung in erzeugten Images bietet Zephyr zudem drei unterschiedliche Targets. Die Werkzeuge laufen demnach auf dem final erzeugten Image und stellen Informationen über die Größe von Symbolen und Code dar, welche sowohl im RAM als auch im ROM verwendet werden. Außerdem kann mithilfe von Compiler Funktionen auch eine Worst-Case-Stack-Nutzungsanalyse erzeugen. Diese Werkzeuge wurden in dieser Arbeit verwendet:

\subsection{ram\_report}

Alle kompiliiertenobjekte und ihre RAM-Nutzung werden in tabellarischer Form aufgelistet. Aus dieser Tabelle kann man Informationen wie Bytes pro Symbol und dem prozentualen Anteil einsehen. Diese Daten werden basierend auf dem Speicherort des Objekts im Dateisystembaum und der Datei gruppiert, die das Symbol enthält. Dies lässt sich mit folgenden West Befehlen umsetzen:

west build -b boardnamefilename

west build -t ram\_report

\subsection{rom\_report}

Alle kompiliiertenobjekte und ihre ROM-Nutzung werden in tabellarischer Form aufgelistet. Aus dieser Tabelle kann man Informationen wie Bytes pro Symbol und dem prozentualen Anteil einsehen. Diese Daten werden basierend auf dem Speicherort des Objekts im Dateisystembaum und der Datei gruppiert, die das Symbol enthält. Dies lässt sich mit folgenden West Befehlen umsetzen:

west build -b boardnamefilename

west build -t rom\_report

\subsection{ram\_plot/ rom\_plot}

Zusätzlich zu den Targets ram\_report und rom\_report erzeugen ram\_plot und rom\_plot Speichernutzungsberichte, aber in Form von Diagrammen. Die Segmente sind interargierbar und hilft dabei, durch die Verzeichnisstruktur zu navigieren und weitere Details auslesen zu können. Sobald der folgende Befehl ausgeführt wird, wird zuerst ein CLI-Bericht erstellt und anschließend ein Browser-Fenster geöffnet:

west build -b boardname filename

west build -t ram\_plotoderrom\_plot

\subsection{Dashboard}

Durch dieses Werkzeug lässt sich ein HTML-Dashboard erzeugen, welches die verschiedenen Werkzeugausgaben und Artefakte in einer einzige Ansicht zusammenfasst. Neben der zusammenfassung der Build-Ergebnisse sind zusätzlich vollständige Speicherberichte zu RAM und ROM in Form von Tabellen und Diagrammen enthalten. Kconfig-Symbolwerte und deren Quekllen werden ebenfalls generiert. Zudem wird ein navigierbarer Devicetree-Ansicht erzeugt, in welchen weitere Details eingesehen werden können. Dies wird durch folgenden Befehl ermöglicht:

west build -b boardnamefilename

west build -t dashboard

\subsection{Flashing}

\section{Zephyr - Grundlagen}

\subsection{Kernel}

\subsection{Threading}

\subsection{Scheduling}

\subsection{Synchronization}

\subsection{Interrupts}

\subsection{Timing}

\subsection{Memory}

\section{Messgrundlagen}

\subsection{Logging}

\subsection{Tracing}

\subsection{Instrumentation}

\subsection{CPU load}

\section{Konfiguration}

\subsection{CMake}

\subsection{Devicetree}

\subsection{Kconfig}

\subsection{Snippets}

\subsection{Flashing}

\subsection{Sysbuild?}

\section{Hardware – Cortex M4F}

\subsection{GPIO}

\subsection{Timer}

\subsection{UART}

\subsection{Bluetooth}

\subsection{SPI/ I2C}

https://www.zephyrproject.org/learn-about/

https://thesis.dial.uclouvain.be/server/api/core/bitstreams/c6169167-70c9-42c0-9654-c01f87638fb0/content

https://link.springer.com/book/10.1007/978-1-4614-3143-5

https://link.springer.com/book/10.1007/978-3-031-28701-5

https://link.springer.com/book/10.1007/978-3-031-73098-6

https://www.intel.com/content/www/us/en/developer/articles/community/zephyr-story-how-became-self-sustaining-ecosystem.html


